\scanpage{179}
\mspage{85}

\uline{Remarque} Lorsque les $D_{d_*}^{*}$ sont tous des $K(\pi,1)$ (ce qui est le cas dans la situation de Teichmüller!) alors la connaissance du système des $\Pi_1\unsure{V}_{d_*,d''_0}^{d'_0}$ précédents (avec les morphismes canoniques entre ces groupoïdes) doit permettre de reconstituer, non seulement le \deleted{type d'homotopie} \inserted{groupoïde} fondamental de $X$ et des \inserted{divers} topos canoniques associés à $X$, mais aussi le topos quasi-galoisien formé des faisceaux sur les topos qui sur chaque topos ``élémentaire'' constitutif sont loc. constants --- et par là aussi les types d'homotopie des topos envisagés --- y compris $X$\ldots

\margin{cette reconstitution doit se faire \unsure{indépendt} de tt autre hypothèse (genre $K(\pi,1)$) de façon essentiellement triviale\ldots}

\uline{Remords} Ce n'est pas le cas pourtant sans hyp. supplémentaire, \unsure{pour} une \unsure{variété}

\margin{178}
